\section{Method}
\label{sec:method}

EvoFSM-RL is defined independently of any particular benchmark: it specifies an
in-context knowledge representation, an online procedure that co-adapts that
representation and the policy weights, and a two-phase schedule that ties
source-pool pretraining to target-app deployment. Section~\ref{sec:exp}
instantiates the same method at two levels of generalization---within a single
benchmark and across two benchmarks---by changing only which apps play the role
of source and target. This section describes the method itself.

\subsection{Setup}
\label{sec:method-setup}

Our setting is the standard mobile GUI agent loop. At each step the agent
receives a screenshot of the current Android screen and a natural-language goal;
it emits one of a fixed set of low-level actions (click on a coordinate or a
labeled element, type a string, scroll, navigate, or signal task completion);
the action is executed on the device, producing the next screenshot. The episode
ends when the agent signals completion, when an external grader verifies
success, or when a per-task step budget is exhausted.

We use Qwen3-VL-8B-Instruct~\citep{qwen3vl2025} as the vision-language backbone
throughout. The surrounding agent harness depends on the benchmark and is
described with each level in Section~\ref{sec:exp} (an action-plus-summary loop
with an accessibility channel for the within-benchmark level, and a pure-vision
native tool-call loop for the cross-benchmark level). The method's contribution
is orthogonal to the harness: it is the structured in-context knowledge $F$ and
the joint adaptation of $F$ together with a LoRA adapter on the policy. The
action-selection prompt has a fixed prefix, an optional slot for $F$, and a
dynamic block containing the goal, history, and current UI. Different baselines
instantiate $F$ differently---empty in the zero-shot baseline, a static $L_C$
text in the static-knowledge baseline, an evolving FSM under our method---but
everything else is byte-identical across baselines, so comparisons isolate the
contribution of $F$ rather than confounding it with prompt differences.

\subsection{Two-layer FSM}
\label{sec:method-fsm}

The core design choice is to structure $F$ as a finite state machine separated
into two layers that differ in their transferability. This separation is the
unit of what gets aggregated, retrieved, and mutated across apps; it is also the
structural commitment that lets the system specialize without breaking transfer.

\textbf{LAYER~1 (app-specific).} A directed graph whose nodes are screens (e.g.\
\texttt{FILE\_LIST}, \texttt{NEW\_FILE\_DIALOG}) and whose edges are interactions
that transition between them. Each node carries visual cues (e.g.\ \emph{``red
plus FAB bottom-right''}), resource hints (the strings the model is likely to
read on that screen), and known dead-ends. Every entry binds to one app's UI and
is therefore non-transferable.

\textbf{LAYER~2 (category-generic).} A set of task-shape archetypes such as
\texttt{ADD\_ENTRY}, \texttt{QUERY\_INFO}, \texttt{REMOVE\_ENTRY}, and
\texttt{FILTER\_LIST}, each with four fields: \emph{precondition},
\emph{abstract\_steps}, \emph{failure\_modes}, and
\emph{verification\_checklist}. LAYER~2 is written without any app-specific
token. A linter (\texttt{lint\_layer2}) rejects any LAYER~2 block, or any
\emph{diff} during online evolution, that contains an app name or a resource
identifier from the source pool. This enforced separation is what makes the
upper layer the structurally guaranteed unit of cross-app transfer.

\textbf{The category-level abstract library $L_C$.} Same-category source apps
share their LAYER~2 content into one merged $L_C$ file per Play Store category.
The merger keeps only claims supported by multiple sources, explicitly
enumerates cross-source failure modes, and never names a specific app. At
deployment, $\mathrm{resolve\_L\_C}(a_{\text{tgt}})$ returns the $L_C$ for the
target app's category if that category is in the source pool, or $\bot$
otherwise. The latter case gives our evaluation a built-in null control: with no
$L_C$ to inject, the prompt degenerates to the zero-shot baseline
byte-identically.

\textbf{How the FSMs are built.} For each source app we collect a handful of
trajectories under the zero-shot baseline, concatenate them into a compressed
text bundle, and ask a frozen Claude Opus model to synthesize a two-layer FSM in
JSON, prompted to keep LAYER~2 strictly app-agnostic; we then lint and audit the
result. The per-app FSMs feed an aggregator that, again with Claude Opus, merges
same-category LAYER~2 blocks into one $L_C$ per category. The concrete source
inventory and the resulting set of category libraries differ by level and are
given in Section~\ref{sec:exp} (for instance, the within-benchmark level
aggregates 12 source apps into six category libraries).

\subsection{The EvoFSM-RL adaptation loop}
\label{sec:method-loop}

At adaptation time we maintain a population $\mathcal{S}_a = \{(F_i, \mu_i,
\sigma_i)\}$ of FSM variants for the current app $a$, each with a TrueSkill
rating~\citep{herbrich2007trueskill}. One iteration proceeds as in
Algorithm~\ref{alg:evofsmrl}.

\begin{table}[t]
\small
\centering
\begin{tabular}{p{0.92\columnwidth}}
\hline
\textbf{Algorithm 1: One EvoFSM-RL iteration} \\
\hline
\textbf{Input:} population $\mathcal{S}_a$, task set $\mathcal{T}_a$, policy
$\pi_\theta$, frozen reference LLM $\pi_{\text{ref}}$, GRPO buffer $\mathcal{B}$ \\
\hline
1. Sample task $t \sim \mathcal{T}_a$ \\
2. Select $M$ variants $\{F_i\}_{i=1}^M$ from the latest-$K$ window of
$\mathcal{S}_a$ with $p_i \propto \exp((\mu_i + \lambda\sigma_i)/T)$ \\
3. For each $i$, run $N$ rollouts of $\pi_\theta$ under prompt $F_i$ on task $t$
with paired seeds; record reward $r_{i,j}$ and per-step log-probs \\
4. Update $\{(\mu_i, \sigma_i)\}$ by TrueSkill from the ranking of $\{V_i =
\tfrac{1}{N}\sum_j r_{i,j}\}$ \\
5. Append the $M{\times}N$ trajectories to $\mathcal{B}$; if $|\mathcal{B}| \ge
K_{\text{buf}}$ run a GRPO step on $\mathcal{B}$ and clear it
(Section~\ref{sec:method-grpo}) \\
6. Let $k = \arg\max_i V_i$. If $r_k > 0$ or $i \bmod 3 = 0$: synthesize a
LAYER~2 diff with $\pi_{\text{ref}}$, apply it to $F_k$, append the child to
$\mathcal{S}_a$ with $\sigma_k + \Delta\sigma$ \\
\hline
\end{tabular}
\caption{One iteration of the EvoFSM-RL adaptation loop. The same procedure runs
in Phase~1 (source-pool pretraining) and Phase~3 (TTA on a target app); the
differences are which $\mathcal{T}_a$ is drawn from and whether $F$ is allowed
to mutate.}
\label{alg:evofsmrl}
\end{table}

Three design choices in this loop are load-bearing.

\textbf{Mutation uses a frozen reference, not the trainable policy.} We prompt
the frozen Claude Opus model with a layered-reflection template over the recent
trajectory window and ask for a JSON-typed diff over LAYER~2 only; the lint pass
guarantees no leakage. We deliberately do not use $\pi_\theta$ itself for the
reflection because, in joint training, the RL-updating policy drifts as it
learns, and using a moving target as a critic of its own behavior degrades the
quality of the proposed edits---a finding ablated in~\citet{zhang2026espl}.

\textbf{Paired seeds within a tournament round.} All $M$ variants are rolled out
on the same task and the same rollout seed within an iteration. The only source
of reward difference between variants is therefore the FSM content, which lets
TrueSkill compare them on identical content rather than confounding the
comparison with task-difficulty noise.

\textbf{LAYER~1 is frozen during online adaptation.} LAYER~1 is built offline
from source-pool trajectories and is not mutated during online adaptation; only
LAYER~2 is. This keeps online evolution bounded to the transferable part of the
knowledge---if LAYER~1 were mutable, a few iterations on a target app could
absorb app-specific edits into the wrong layer and silently break transfer
downstream.

\subsection{Group-relative policy optimization}
\label{sec:method-grpo}

The weight update is a group-relative policy
gradient~\citep{shao2024deepseekmath} computed on the buffered trajectories. We
compute the reward for each rollout as
\begin{equation}
R(t,\tau) = s + w_e \cdot \max\!\left(0,\, 1 - \tfrac{n_{\text{steps}}(\tau)}{B(t)}\right),
\label{eq:reward}
\end{equation}
where $s \in \{0, 0.5, 1\}$ is the grader's verdict, $B(t)$ is the per-template
step budget, and $w_e = 0.2$ keeps the efficiency bonus small enough that
success dominates.

Following the single-task M-prompts-by-N-rollouts shape of Algorithm~1
in~\citet{zhang2026espl}, we key groups by the $(F_i, t)$ pair: within each
iteration, the $N$ rollouts under variant $F_i$ on task $t$ form one group, and
the advantage of rollout $j$ is
\begin{equation}
A_{i,j} = r_{i,j} - V_i, \qquad V_i = \tfrac{1}{N}\sum_{k=1}^{N} r_{i,k}.
\label{eq:advantage}
\end{equation}
The signal $A_{i,j}$ is purely a within-(FSM, task) deviation from the local
mean. A group of size~1 has $A \equiv 0$ and contributes nothing to the update,
which is why we require $N \ge 2$ whenever the weight channel is active.

The LoRA-only gradient on a buffer of active trajectories is
\begin{equation}
\nabla_\theta \mathcal{L} = -\sum_{j \in \mathcal{B}_{\text{active}}} \frac{A_j}{T_j} \sum_{t=1}^{T_j} \nabla_\theta \log \pi_\theta(a_t \mid s_t, F),
\label{eq:grpo-loss}
\end{equation}
where $T_j$ is the number of agent steps in rollout $j$. The $1/T_j$
normalization keeps long and short rollouts on the same gradient scale; without
it, the raw gradient norm scales with trajectory length and gets clipped
uniformly, eliminating the actual learning direction in favor of a
trajectory-length artifact. We freeze the base vision-language model and update
only LoRA matrices on the attention projections $\{q_{\text{proj}},
v_{\text{proj}}\}$~\citep{hu2021lora}, optionally anchored to the frozen
reference policy with a small KL penalty (per-level coefficients in the
appendix).

\subsection{Two-phase application}
\label{sec:method-phases}

The same EvoFSM-RL loop runs at two points in the pipeline. Using the
\emph{same} loop in both places---rather than training one model offline and
fine-tuning another online---is what lets us call the deployment-time procedure
a well-defined operation rather than an ad-hoc fine-tune.

\textbf{Phase~1: shared-LoRA pretraining.} A single training run over the source
pool produces a shared LoRA adapter $\pi^{\text{pre}}_\theta$. Each iteration
samples one source app uniformly, samples one of its tasks, runs $N$ rollouts
under the app's static FSM, and accumulates the trajectories into a single GRPO
buffer. The buffer's groups are still keyed by $(F_i, t)$, so gradients from
rollouts on different source apps flow through their respective task groups into
the same adapter. FSM mutation is frozen in Phase~1: the static FSMs and $L_C$
are produced by simpler non-RL pipelines and serve as inputs to this phase.

\textbf{Phase~3: deployment-time adaptation.} For each target app we run the same
loop on the app's $T_{\text{adapt}}$ split, starting from
$\pi^{\text{pre}}_\theta$ and seeding the population from the category-matched
$L_C$. Both the FSM population and the LoRA adapter evolve online. After a fixed
budget of adaptation iterations, the LoRA adapter and the FSM champion are
frozen and the agent is evaluated on the held-out $T_{\text{eval}}$ split. This
is the test-time half of the framework, and what we call B4 in the ablation; B3
is the same loop with the weight channel disabled.

\textbf{Instantiation at each level.} Phase~1 and Phase~3 differ between the two
evaluation levels only in which apps fill the source and target roles. At the
within-benchmark level, the source pool and the target apps are disjoint app
sets drawn from the \emph{same} benchmark, and adaptation/evaluation are
template-disjoint within each target app. At the cross-benchmark level, Phase~1
pretrains on the \emph{entire} training benchmark and Phase~3 adapts on each
target app of a \emph{separate} deployment benchmark, with adaptation and
evaluation task-disjoint within that benchmark. Section~\ref{sec:exp} gives the
concrete app inventories and split sizes.

A note on what Phase~2 looks like in our implementation. The full algorithmic
description specifies a target-app cold-start that uses $\pi^{\text{pre}}_\theta$
to roll out a few exploration trajectories and synthesize a target-app LAYER~1
before adaptation begins. We collapse this to a single category look-up: $L_C
\leftarrow \mathrm{resolve\_L\_C}(a_{\text{tgt}})$. The target app's LAYER~1 is
rediscovered in-context during the B3/B4 loop itself rather than offline. This
works because the visual encoder is strong; we discuss when it fails in
Section~\ref{sec:limitations}.

The Phase~1 $\to$ Phase~3 split is what makes the joint paradigm tractable at the
deployment-time budget. A small number of iterations from a randomly initialized
LoRA is not enough to learn anything meaningful, whereas the same budget starting
from $\pi^{\text{pre}}_\theta$ exercises the policy in a regime where the local
gradient signal is small but the starting point already encodes a useful prior.
